\chapter{误差、理论保证与可核对的推导实例}
\label{th:guarantees}

\section{生成误差至少有五个来源}

一个可追溯的链路应分开记录：表示误差、先验/端点误差、
场或映射的学习误差、数值离散误差，以及引导/条件近似误差。
它们不能简单用一个训练 MSE 相加，但分开定义可避免把任一失败都归因于模型大小。

表示误差来自 codec 或离散化，先验误差来自实际 $\pi_T\ne p_T$。
学习误差指真实边缘对象与网络输出的偏差。
数值误差是在已学对象上执行有限步的偏差，
引导误差还可能改变理论目标本身。
蒸馏把一部分积分成本转为映射拟合误差，不能据“少了求解器”认定总误差为零。

\section{ODE 稳定性：局部误差如何传播}

设真实速度 $u$ 关于状态为 $L$-Lipschitz，
在所比较轨迹区域内 $\|u_\theta-u\|\leq\delta(t)$。
把两条轨迹初值误差记为 $e_0$，Gronwall 不等式给出
\begin{equation}\label{th:gronwall}
\|x_t^\theta-x_t\|\leq e^{Lt}e_0+
 \int_0^t e^{L(t-r)}\delta(r)\,\mathrm dr.
\end{equation}
反向生成可用反时钟得到同形式的界。
若再加入数值离散误差，会多出依赖步长和高阶导数的项。
强引导可能增大有效 Lipschitz 常数，
高阶公式在非光滑裁剪与奇异端点附近也可能失去名义收敛阶。

这是沿指定区域的一致场误差结论。
训练得到的 $\E_{p_t}\|u_\theta-u\|^2$ 小，
不等于已经证明上面使用的全区域上界。
分布外误差、长尾与低密度通道需要额外控制。

\section{SDE 的相对熵界与 ODE 界有不同前提}

式\eqref{th:path-kl} 将时间加权分数误差联系到反向 SDE 路径 KL。
借数据处理可限制端点 KL；Pinsker 不等式又给出
$\mathrm{TV}(P,Q)\leq\sqrt{\KL(P\Vert Q)/2}$。
这是一条有假设的逻辑链，不是无条件的模型质量保证。

两过程扩散不同、端点支持不相容、相关积分发散时，该等式不能原样使用。
确定性 ODE 轨迹测度可能相互奇异，不能照抄 SDE 的路径 KL 证明。
实际 DDPM、SDE 数值离散与可训练流映射各有自己的误差分析。

终端先验近似、$t=\varepsilon$ 截断、密度光滑性、尾部矩与网络正则性
往往出现在定理条件中。引用一个收敛定理时，必须连同这些条件和误差度量一起陈述，
而不是只写“已证明能收敛到数据”。

\section{高维奇异端点与数据流形}

干净数据可能集中在低维结构上，$p_0$ 未必具有全空间的平滑密度。
正噪声 $b_t>0$ 通常使高斯平滑后的密度良好，
但当 $t\downarrow0$ 时分数与高阶导数可能发散。
从0积分到1的形式公式不意味着每个端点都可直接执行。

常见做法是停在 $\varepsilon>0$，再执行指定去噪/解码步骤，
并单独报告其误差。
后验均值可能带来细节平滑，额外随机采样可能改变方差；
应按生成模型选择，不能把最后一步当任意美化插件。

\section{似然、FID、条件遵循与覆盖度分别测什么}

似然评估分布赋给真实样本的密度，受去量化、codec 与数值估计口径影响。
FID 在特征空间用均值与协方差比较高斯近似，
不同分布可以具有相同前两阶矩，因此低 FID 不证明分布完全一致。
估计样本数、特征实现和预处理也影响结果。

条件一致性、编辑保持、文字精度、多样性与校准是其它维度。
用少量挑选样本证明理论正确，或用单一奖励证明整体质量，都会混淆层次。
训练集重叠与重复素材会使质量估计失真；
少步模型还需在相同 NFE、引导与预算口径下比较。

\section{五个可以自行核对的推导}

\subsection{高斯分数与去噪后验}

设 $X_0\sim\mathcal N(m,S)$，$X_t=aX_0+b\epsilon$。
则 $C=a^2S+b^2I$，
\[
s_t(x)=-C^{-1}(x-am),\qquad
\E[X_0\mid x]=m+aS C^{-1}(x-am).
\]
把右侧代入 $a\E[X_0\mid x]=x+b^2s_t(x)$ 即可核对 Tweedie。
在标量 $m=0,S=1,a=0.8,b=0.6$ 时，$C=1$、$D^*=0.8x$。
后验方差仍为 $0.36$，所以用 $0.8x$ 代替后验随机样本会丢失不确定性。

\subsection{DDPM 后验由高斯精度相加得到}

把与 $x_{k-1}$ 有关的指数写成
\[
\frac{\|x_k-\sqrt{\alpha_k}x_{k-1}\|^2}{\beta_k}
 +\frac{\|x_{k-1}-\sqrt{\bar\alpha_{k-1}}x_0\|^2}{1-\bar\alpha_{k-1}}.
\]
二次系数为 $\alpha_k/\beta_k+1/(1-\bar\alpha_{k-1})$
$=(1-\bar\alpha_k)/(\beta_k(1-\bar\alpha_{k-1}))$。
取倒数得到 $\tilde\beta_k$，线性系数乘这个方差得到第\ref{th:ddpm}章的均值。
这也说明原始单步方差与后验方差不是同一个量。

\subsection{一步流匹配的反例}

第\ref{th:flow}章独立标准高斯的插值有 $c_t=(1-t)^2+t^2$，
$\dot c_t=4t-2$。线性 ODE $\dot x=\dot c_t x/(2c_t)$ 的精确解
$x_t=\sqrt{c_t/c_s}\,x_s$，两端 $c_0=c_1=1$，故端点映射为恒等。
一步 Euler 却把噪声压到0，展示条件直线监督与边缘积分的区别。

\subsection{MeanFlow 恒等式的解析验证}

对标量速度 $u(x,t)=kx$，
$\Phi_{s\to t}(x)=e^{k(t-s)}x$，所以
\[
U(x,s,t)=\frac{e^{k(t-s)}-1}{t-s}x.
\]
代入式\eqref{th:meanflow-identity}，对 $s$ 求导时目的时刻 $t$ 固定，
可验证两边相等。若遗漏 $D_xU\,u$ 或把 $s$ 的导数符号写反，等式立刻失效。
在 $t\to s$ 时应用极限 $U\to kx$，不直接计算 $0/0$。

\subsection{掩码反向概率的归一与方向}

令 $a(s)=0.8,a(t)=0.5$。
已掩码 token 在较早时刻为干净状态的概率是 $0.3/0.5=0.6$，
仍被掩码的概率是 $0.2/0.5=0.4$。
若把 $a(s)-a(t)$ 写成相反顺序，就得到负概率，
这比只看一个“反向公式”更容易发现时钟混用。

\Needspace{20\baselineskip}
\section{本次补全的完成边界}

本篇补齐了主参考的三种起源、统一视角、OT、引导、采样、蒸馏和流映射，
并加入 v3 的离散概率演化。
每章是按本书符号重新组织的中文讲解，不是参考书逐页翻译。
具体理论分析范围和下载身份记录在随附的来源清单；
未复制原书全文，未运行模型训练或把数学检查称为模型复现。

读者继续研究时，可优先追问：路径是否固定、监督是否无偏、边界是否可逆、
近似是否改变目标、误差在哪个分布下度量。
这些问题比记住某个方法名称更能帮助判断后续新工作的实际贡献。

\begin{theorytakeaway}
理论保证要带着假设、误差度量与端点一起读。
解析恒等式、数值核对、训练实验和产品验收提供不同证据，不能互相替代。
\end{theorytakeaway}
\threading{主参考第6--7、11--12章及理论附录\cite{theorybook}，
连续误差框架联系\cite{sde}，流映射与离散比值分别联系\cite{meanflow,sedd}。}
